\documentclass[a4paper,11pt]{article}

% Packages
\usepackage[utf8]{inputenc}
\usepackage[T1]{fontenc}
\usepackage{amsmath, amssymb}
\usepackage{geometry}
\usepackage{xcolor}
\usepackage{enumitem}
\usepackage{fancyhdr}
\usepackage{array}
\usepackage{booktabs}
\usepackage{graphicx}
\usepackage{titlesec}
\usepackage[
    colorlinks=true,
    linkcolor=black,
    urlcolor=black
]{hyperref}

% Page setup
\geometry{margin=2.5cm}
\setlength{\headheight}{14pt}
\pagestyle{fancy}
\fancyhf{}
\fancyhead[L]{\textit{Intelligent Assistant for L1 Courses}}
\fancyhead[R]{\textit{CMI Informatique --- 2025/2026}}
\fancyfoot[C]{\thepage}

% Style of sections
\titleformat{\section}
    {\large\bfseries}
    {\thesection.}{0.6em}{}
    [\vspace{2pt}\titlerule\vspace{4pt}]

\titleformat{\subsection}
    {\normalsize\bfseries}
    {\thesubsection.}{0.6em}{}

% ============================================================
\begin{document}

% ============================================================
% COVER OF RESEARCH PROJECT
% ============================================================
\thispagestyle{empty}

\vspace*{3cm}

\begin{center}
    {\small\scshape Research Project --- CMI Informatique}\\[1cm]
    {\rule{\linewidth}{0.4pt}}\\[0.6cm]
    {\LARGE\bfseries Creation of an Intelligent Assistant\\[0.3cm]for First-Year Courses}\\[0.6cm]
    {\rule{\linewidth}{0.4pt}}\\[1.2cm]
    {\normalsize\itshape Target audience: Everyone}\\[0.4cm]
    {\normalsize Supervisor: Yannick Estève}\\[2.5cm]
    {\normalsize
        A. Louvet \quad L. Ouelhadj \quad T. De Faveri-Serre \quad C. Cloupet \quad H. Baudey
    }\\[0.8cm]
    {\small Academics Years 2025--2027}
\end{center}

\vfill

\begin{center}
\noindent\rule{0.5\linewidth}{0.4pt}\\[0.5cm]
{\small
This project aims to one one hand make  some research contributions about
LLMs, RAG, vector search and explain it regardless of prior familiarity 
with the concept. \\ and orchestration frameworks in the context of 
educational assistants, in the order build a conversational AI assistant grounded in 
Retrieval-Augmented Generation (RAG), capable of answering students'
questions based on their own course documents.
}\\[0.4cm]
\noindent\rule{0.5\linewidth}{0.4pt}
\end{center}

\vspace{2cm}

% ============================================================
% TABLE OF CONTENTS
% ============================================================

\newpage
\tableofcontents

% ============================================================
% PART 1: STATE OF THE ART (DEFINITIONS AND CONTEXT)
% ============================================================

\newpage
\section{State of the Art : Definitions and Context}

\subsection{Setting the Stage with Artificial Intelligence}

\subsubsection{A Brief History}
The idea of machines that can think is not new. As early as 1950, Alan
Turing proposed his famous \emph{Turing Test} as a criterion for machine
intelligence \cite{turing1950computing}: if a machine can sustain a
conversation indistinguishable from a human, it could, in some sense, be
said to think. The term ``Artificial Intelligence'' itself was coined in 1956
by John McCarthy at the Dartmouth Conference, which is widely regarded as
the founding moment of the field.
 
The decades that followed alternated between periods of optimism and
so-called ``AI winters'' --- intervals of reduced funding and interest
caused by the gap between ambitious promises and modest results. The
landscape shifted dramatically in the 1980s and 1990s with the rise of
\emph{machine learning}: rather than hard-coding rules into a program,
researchers began designing systems that could \emph{learn} patterns
directly from data. This paradigm shift laid the foundations for everything
that followed.

The true turning point came in 2012, when a deep neural network called
AlexNet \cite{krizhevsky2012alexnet} crushed the competition in the ImageNet
visual recognition challenge, demonstrating that \emph{deep learning} ---
neural networks with many layers trained on large datasets --- could achieve
performance far beyond classical methods. From that moment, deep learning
became the dominant approach across almost every subfield of AI.

\subsubsection{So What Is Artificial Intelligence?}
The term \emph{Artificial Intelligence} is one of the most used --- and most
misused --- expressions in today's technological landscape. Before discussing
the specific systems at the heart of this project, it is worth taking a step
back to understand what AI actually is, where it comes from, and why it has
become so central to modern computing.

At its core, Artificial Intelligence refers to any computational system
designed to perform tasks that would normally require human intelligence ---
such as recognizing objects in an image, understanding a spoken sentence,
or deciding the next move in a chess game. A more precise definition,
proposed by Russell and Norvig in their seminal textbook \cite{russell2010ai},
describes an AI agent as a system that \emph{perceives its environment through
sensors and acts upon it through actuators in order to maximize a performance
measure}.
 
A useful distinction is that between \emph{narrow AI} (also called weak AI)
and \emph{general AI} (or strong AI). Narrow AI --- which covers virtually
all systems in existence today --- is designed to excel at one specific type
of task: a spam filter classifies emails, a recommendation engine suggests
movies, a language model generates text. General AI, by contrast, would be
a system capable of reasoning across arbitrary domains at a human level; it
remains largely hypothetical and is an active subject of long-term research.
Every AI system discussed in this project falls squarely in the narrow AI
category.
% ------------------------------------------------------------
\subsection{A closer look at Large Language Models (LLMs)}
ceceeeedvfieij (à refaire)
\subsubsection{Limitations: Hallucinations and Knowledge Boundaries}
 
Despite their impressive performance, LLMs suffer from several
well-documented limitations. The most critical one in an educational context
is the tendency to produce \emph{hallucinations} --- generating text that is
fluent and plausible but factually incorrect or unfounded
\cite{zhang2023hallucination, huang2025hallucination}.
 
Hallucinations arise from multiple causes. First, LLMs rely on statistical
correlations in their training data rather than grounded factual reasoning.
Second, their knowledge is static: they have a training cutoff date and
cannot access new or domain-specific information unless explicitly provided.
Third, when a query falls outside the model's confident knowledge region, it
may confabulate an answer rather than acknowledging uncertainty. According to
a 2024 survey published in the \emph{ACM Transactions on Information
Systems}, hallucinations are categorized into two main types:
\emph{factuality hallucination} (stating incorrect facts) and
\emph{faithfulness hallucination} (producing content inconsistent with a
given source).
 
These limitations are particularly problematic in an educational context,
where accuracy is essential. An assistant that confidently provides wrong
information about a course concept could actively mislead students. This
motivates the use of external document grounding through RAG.

\subsubsection{Biase Considerations Issue : Feedback Loop Degradation via RLHF}
A subtler but  significant limitation concerns the feedback mechanism used to
improve LLMs by using techniques such as \emph{Reinforcement Learning from Human Feedback} 
(RLHF) and Constitutional AI are used to align model behaviorwith human preferences 
and reduce harmful outputs. This system bases on a system in which users
rate model responses positively or negatively and the model is updated
accordingly \cite{ouyang2022rlhf}. While this process is intended to align
model behavior with user expectations, it can produce perverse effects when
users consistently submit poorly formulated prompts.
 
When a user provides an ambiguous, incomplete, or incoherent query, the model
has no reliable basis on which to produce a relevant answer. Then it generates a
response that is statistically plausible given the input, but which
inevitably fails to match the user's unstated intent. The user, perceiving
the response as irrelevant or absurd, assigns a negative rating --- even
though, from the model's perspective, the response was a reasonable
completion of the given prompt. This misalignment between the signal (a
negative rating) and its actual cause (a poorly formulated query rather than
a model error) introduces \emph{corrupted supervision} into the training
loop.
 
Repeated at scale, this dynamic causes the model to associate certain
legitimate response patterns with negative reward and to unlearn them,
leading to a measurable degradation in general response quality over time. HAS
study by Chen et al. \cite{chen2023chatgpt} empirically documented this
phenomenon for ChatGPT, observing significant behavioral shifts between
model versions, with certain tasks where earlier versions outperformed later
ones. The authors attribute part of this regression to the compounding
effects of user feedback noise in the RLHF pipeline.
 
This issue is particularly relevant in an educational setting, where students
--- especially at the first-year level --- may lack the prompt-engineering
literacy to formulate precise queries. An assistant that degrades under
noisy feedback could thus become progressively less useful for its primary
audience, underscoring the importance of designing robust feedback collection
mechanisms and, where possible, reducing reliance on unfiltered user ratings
for model updates.
% ------------------------------------------------------------
\subsection{Let's recap}
A natural starting point for this overview is a question that seems simple
but is actually quite nuanced: what exactly is an LLM, and how does it
differ from the broader notion of ``AI''? The term \emph{Artificial
Intelligence} (AI) is an umbrella concept encompassing any technique or
system that enables machines to mimic human intelligence, including machine
learning, computer vision, robotics, and natural language processing. A
\emph{Large Language Model} (LLM), by contrast, refers to a specific
category of AI model designed to understand and generate human language.
LLMs are a subset of AI, just like \emph{Machine Learning } (ML) or \emph{Deep Learning} (DL) etc\dots
while all LLMs are AI systems, the converse is not true.
 
Concretely, LLMs are neural networks trained on massive text corpora using
self-supervised learning, based on the Transformer architecture introduced
by Vaswani et al.\ \cite{vaswani2017attention}, which enables models to
capture long-range dependencies in text through self-attention mechanisms.
The subsequent introduction of OpenAI's GPT series --- GPT-2 in 2019,
GPT-3 in 2020, and later GPT-4 --- demonstrated that scaling up model size
and training data leads to \emph{emergent capabilities} in language
understanding and generation \cite{wei2022emergent}: abilities that appear
qualitatively only beyond a certain scale threshold. Contemporary LLMs such
as GPT-4, Anthropic's Claude, Meta's LLaMA, Google's Gemini, and Mistral
exhibit remarkable proficiency across a wide range of tasks --- from question
answering and summarization to code generation and multi-step reasoning ---
and are characterized by parameter counts in the tens or hundreds of billions.
These models operate by encoding vast amounts of world knowledge into their
parameters through statistical pattern learning over billions of tokens, and
their trajectory of development has been defined by increasing model size,
diversifying training data, and refining training techniques.
% ------------------------------------------------------------
% ------------------------------------------------------------
\subsection{Retrieval-Augmented Generation (RAG)}
 
\subsubsection{Principle and Motivation}
Before diving into the technical architecture, it is worth stepping back to
explain what RAG is and why it was developed, so as to make this research
accessible and interesting to everyone.
 
Imagine asking a very knowledgeable friend a question about a specific topic
--- say, the content of a course they have never attended. No matter how
smart they are, they simply cannot know what was taught in that particular
class. An LLM faces the exact same limitation: it only knows what it saw
during training, and it has no way to consult documents it has never
encountered. Now imagine instead that before answering, your friend is
allowed to quickly flip through the relevant course notes and use what they
find to inform their reply. That, in essence, is what RAG does: it gives the
model access to a specific set of documents at the moment of answering, so
that its response is grounded in those sources rather than in its general
(and potentially outdated or incomplete) memory.
 
More formally, Retrieval-Augmented Generation (RAG) is a paradigm that
enhances LLMs by connecting them to an external knowledge base at inference
time \cite{lewis2020rag}. Rather than relying solely on the model's
parametric memory, a RAG system first retrieves relevant document passages
from a curated corpus based on the user's query, then passes these passages
as context to the LLM, which generates a grounded response. 

This approach directly addresses the limitations described in the previous
section: it provides the LLM with up-to-date and domain-specific
information, and it anchors the generation in explicit, verifiable sources
--- significantly reducing the risk of hallucination. A survey specifically
focused on RAG for educational applications \cite{rag_education_survey}
confirms that this approach substantially improves the relevance and factual
accuracy of AI-generated answers in academic settings.

\newpage
\subsubsection{RAG Architecture}
 
A standard RAG pipeline consists of three stages: 
 
\begin{enumerate}
  \item \textbf{Indexing (developer phase):} Source documents (PDFs, course
    slides, lecture notes, etc.) are split into chunks, encoded as vector
    embeddings using an embedding model, and stored in a vector database.
    \emph{(cf Section 1.4 with details about vector search and indexing)}.

  \item \textbf{Retrieval (user phase):} When a user submits a query, it is
    encoded into the same embedding space. The system then performs a
    similarity search to retrieve the top-$k$ most relevant chunks from the
    database.
 
  \item \textbf{Generation:} The retrieved chunks are prepended to the user
    query as context in the LLM's prompt. The LLM then generates a response
    grounded in this retrieved information.
\end{enumerate}

\subsubsection{RAG Variants and Evolution}
 
The field has matured significantly since the original RAG paper. Surveys
from 2024--2025 identify three main paradigms \cite{gao2024rag}:
 
\begin{itemize}
  \item \textbf{Naive RAG:} The basic retrieve-then-generate pipeline
    described above.
  \item \textbf{Advanced RAG:} Introduces query rewriting, re-ranking of
    retrieved results, and iterative retrieval to improve precision.
  \item \textbf{Modular RAG:} A more flexible architecture where each
    component (retriever, reader, memory, etc.) can be independently swapped
    or optimized.
\end{itemize}
 
More recent work also explores \textbf{Agentic RAG}, where the model can
decide \emph{when} to retrieve and \emph{which} sources to consult based on
the complexity of the query, using a multi-step reasoning loop
\cite{zhao2024rag}.
 
For our project, a Naive or Advanced RAG architecture is sufficient given
the well-defined and bounded scope of the document corpus (first-year course
materials).
% ------------------------------------------------------------
% ------------------------------------------------------------
\subsection{Vector Search and Indexing}

% ============================================================
% PART 1.2: PROJECT OVERVIEW AND SEARCHES ABOUT AI ASSISTANTS IN EDUCATION
% ============================================================
\newpage
\section{Project Overview}

\subsection{Pedagogical Objectives}
The goal is to discover the fundamental principles of modern artificial intelligence through
large language models (LLM) and the RAG (Retrieval-Augmented Generation) system.
Students will learn how to interface an LLM with a custom document base to create
a conversational assistant capable of answering questions about their courses.
\subsection{AI Assistants in Educational Settings}
 
\subsubsection{Overview of Existing Work}
 
The application of AI-powered assistants to education has attracted growing
research interest. A survey published in \emph{Applied Sciences}
\cite{mdpi2025rag} identified 47 studies on RAG-based educational chatbots,
the majority published in 2024, demonstrating the rapid growth of this field.
 
Several real-world deployments are directly relevant to our project:
 
\begin{itemize}
  \item \textbf{UniMiB Assistant} (Bicocca University, Milan, 2024)
    \cite{unimib2024}: A RAG-based student assistant capable of answering
    questions about university procedures, course schedules, and
    administrative information. It indexes official university documents and
    uses an LLM to generate answers grounded in those sources.
 
  \item \textbf{OwlMentor} (Frontiers in Psychology, 2024)
    \cite{owlmentor2024}: An AI-powered learning environment designed to help
    university students comprehend scientific texts. Integrated into a course,
    it offered document-based chat, automatic question generation, and quiz
    creation.
 
  \item \textbf{CS50 AI Assistant} (Harvard University): Integrated into an
    introductory computer science course, a RAG-based assistant was shown to
    provide detailed and contextually appropriate responses, facilitating
    personalized tutoring at scale.
 
  \item \textbf{Medical education RAG assistant} (Nature \emph{npj Digital
    Medicine}, 2025) \cite{npj2025}: A teaching assistant deployed in a
    medical school course, where usage patterns and student feedback were
    analyzed across two consecutive cohorts, demonstrating the feasibility and
    effectiveness of RAG-based assistants in high-stakes educational contexts.
\end{itemize}
 
\subsubsection{Effectiveness and Challenges}
 
The general finding across these studies is that RAG-based assistants are
effective at providing accurate, context-grounded answers to domain-specific
questions, and that students perceive them positively as a complement to human
instruction. However, several challenges are consistently reported:
 
\begin{itemize}
  \item \textbf{Document quality and coverage:} The quality of the assistant's
    answers is directly bounded by the quality and coverage of the indexed
    document corpus. Poorly formatted or incomplete course documents degrade
    performance.
  \item \textbf{Out-of-scope queries:} Students often ask questions outside
    the domain of the indexed documents. Robust systems must detect such
    queries and gracefully decline to speculate.
  \item \textbf{Evaluation difficulty:} Assessing the correctness of generated
    answers in open-ended educational contexts is non-trivial and often
    requires human expert evaluation.
  \item \textbf{Student trust calibration:} Students may over-rely on AI
    answers or, conversely, distrust them even when correct. Communicating
    uncertainty and citing sources is important for appropriate trust
    calibration.
\end{itemize}
 
% ------------------------------------------------------------
\newpage
\subsection{Summary and Positioning of Our Project}
 
This state of the art reveals that the combination of LLMs and RAG represents
the current best practice for building domain-specific conversational
assistants. The technology stack we plan to use --- Python, an embedding model
(\texttt{sentence-transformers} or OpenAI), a vector database (FAISS or
Chroma), and an orchestration framework (LangChain or LlamaIndex) --- is
well-validated both in the academic literature and in real-world educational
deployments.
 
The main originality of our project lies in its specific scope: building an
assistant tailored to \emph{first-year CMI Informatique course materials},
which have specific characteristics (introductory-level computer science,
mathematics, and project content in French and English). The challenges we
anticipate align with those reported in the literature: ensuring sufficient
document coverage, managing out-of-scope questions, and calibrating student
trust in the system's outputs.
\newpage


% Parler de comment fonctionne l'IA (vectorisation, embeddings, etc.)
% Parler de FAISS, Chroma, etc.
% Parler de LangChain, LlamaIndex, etc.
% ============================================================
%  Bibliography entries (to paste into your .bib file)
% ============================================================
%
% @inproceedings{vaswani2017attention,
%   author    = {Vaswani, Ashish and others},
%   title     = {Attention Is All You Need},
%   booktitle = {Advances in Neural Information Processing Systems (NeurIPS)},
%   year      = {2017}
% }
%
% @article{wei2022emergent,
%   author  = {Wei, Jason and others},
%   title   = {Emergent Abilities of Large Language Models},
%   journal = {Transactions on Machine Learning Research},
%   year    = {2022}
% }
%
% @article{zhang2023hallucination,
%   author  = {Zhang, Yue and others},
%   title   = {Siren's Song in the {AI} Ocean: A Survey on Hallucination in
%              Large Language Models},
%   journal = {arXiv preprint arXiv:2309.01219},
%   year    = {2023}
% }
%
% @article{huang2025hallucination,
%   author  = {Huang, Lei and others},
%   title   = {Large Language Models Hallucination: A Comprehensive Survey},
%   journal = {arXiv preprint arXiv:2510.06265},
%   year    = {2025}
% }
%
% @inproceedings{lewis2020rag,
%   author    = {Lewis, Patrick and others},
%   title     = {Retrieval-Augmented Generation for Knowledge-Intensive {NLP} Tasks},
%   booktitle = {Advances in Neural Information Processing Systems (NeurIPS)},
%   year      = {2020}
% }
%
% @article{gao2024rag,
%   author  = {Gao, Yunfan and others},
%   title   = {Retrieval-Augmented Generation for Large Language Models: A Survey},
%   journal = {arXiv preprint arXiv:2312.10997},
%   year    = {2024}
% }
%
% @article{zhao2024rag,
%   author  = {Zhao, Penghao and others},
%   title   = {A Comprehensive Survey of Retrieval-Augmented Generation ({RAG}):
%              Evolution, Current Landscape and Future Directions},
%   journal = {arXiv preprint arXiv:2410.12837},
%   year    = {2024}
% }
%
% @article{rag_education_survey,
%   title   = {Retrieval-Augmented Generation ({RAG}) Chatbots for Education:
%              A Survey of Applications},
%   journal = {Computers and Education: Artificial Intelligence},
%   year    = {2024}
% }
%
% @article{mdpi2025rag,
%   title   = {Retrieval-Augmented Generation ({RAG}) Chatbots for Education:
%              A Survey of Applications},
%   journal = {Applied Sciences (MDPI)},
%   year    = {2025}
% }
%
% @article{owlmentor2024,
%   title   = {Exploring Generative {AI} in Higher Education: A {RAG} System
%              to Enhance Student Engagement with Scientific Literature},
%   journal = {Frontiers in Psychology},
%   year    = {2024}
% }
%
% @article{unimib2024,
%   author  = {Aiello, Luca Maria and others},
%   title   = {{UniMiB} Assistant: Designing a Student-Friendly {RAG}-Based Chatbot},
%   journal = {arXiv preprint arXiv:2411.19554},
%   year    = {2024}
% }
%
% @article{npj2025,
%   title   = {A Generative {AI} Teaching Assistant for Personalized Learning
%              in Medical Education},
%   journal = {npj Digital Medicine},
%   year    = {2025}
% }
%
% @online{langcopilot2025,
%   title   = {Best {RAG} Frameworks 2025: {LangChain} vs {LlamaIndex} vs {Haystack}},
%   url     = {https://langcopilot.com/posts/2025-09-18-top-rag-frameworks-2024-complete-guide},
%   year    = {2025}
% }
%
% @online{ibm2024llamaindex,
%   author  = {{IBM}},
%   title   = {{LlamaIndex} vs {LangChain}: What's the Difference?},
%   url     = {https://www.ibm.com/think/topics/llamaindex-vs-langchain},
%   year    = {2024}
% }
\bibliographystyle{plain}
\bibliography{references}
\end{document}